\documentclass[conference]{IEEEtran}
\IEEEoverridecommandlockouts

\usepackage{cite}
\usepackage{amsmath,amssymb,amsfonts}
\usepackage{graphicx}
\usepackage{textcomp}
\usepackage{xcolor}
\usepackage{booktabs}
\usepackage{url}
\newtheorem{proposition}{Proposition}

\begin{document}

\title{APEX: Latent-Space Kernel Fusion for Selective\\ Prediction with Frozen Medical Image Classifiers}

\author{
\IEEEauthorblockN{Mohammad Kamal Abdulaziz}
\IEEEauthorblockA{\textit{King Abdullah II School of Information Technology}\\
\textit{The University of Jordan}\\
Amman, Jordan \\
moh0217526@ju.edu.jo}
\and
\IEEEauthorblockN{Rizik Al-Sayyed}
\IEEEauthorblockA{\textit{King Abdullah II School of Information Technology}\\
\textit{The University of Jordan}\\
Amman, Jordan \\
r.alsayyed@ju.edu.jo}
}

\maketitle

\begin{abstract}
A diagnostic model allowed to abstain is judged less by how often it is right than
by how well its confidence orders its own mistakes. We present APEX, a post-hoc
layer that improves that ordering without touching the network. APEX leaves the
backbone frozen and works on the penultimate representation the classifier
computes and then discards: it concatenates that latent vector with image
descriptors, projects the pair onto a shared low-dimensional subspace, evaluates
three kernel experts there, and mixes their posteriors with the original softmax
under a confidence-conditional gate. The weights are fitted by minimising negative
log-likelihood, a strictly proper rule, and we show that choice is what keeps the
fused probabilities honest. Across six imaging modalities, three backbone families
and five-fold cross-validation (90 trained heads, 18 paired configurations,
fifteen metrics), APEX lowers the area under the risk--coverage curve by 44.3\%
and selective risk at 90\% coverage from 10.71\% to 6.52\%. AURC, Brier score and
log-likelihood improve in all 18 configurations against each of six baselines at
$p<10^{-4}$. Expected calibration error does not, and that asymmetry is the point:
selective risk is invariant to monotone rescaling of confidence where calibration
error is not, and in 15 of 18 configurations the lowest-calibration-error method is
not the lowest-selective-risk one.
\end{abstract}

\begin{IEEEkeywords}
selective prediction, confidence calibration, post-hoc methods, medical image
analysis, kernel methods, proper scoring rules
\end{IEEEkeywords}

\section{Introduction}

A deployed diagnostic model is rarely asked to decide everything. The realistic
workflow is triage: the system handles the cases it is sure about and refers the
rest to a clinician~\cite{kompa2021second}. The decisive question is then not
aggregate accuracy but whether the model's confidence separates the cases it will
get right from the ones it will get wrong---a property called selective
prediction, summarised by the area under the risk--coverage
curve~\cite{elyaniv2010foundations,geifman2017selective}.

Post-hoc calibration is the natural place to look for it: a calibrator adjusts a
trained model's probabilities on a small held-out split, costs minutes rather than
GPU-days, and leaves the network untouched---which matters when the network is a
frozen foundation backbone or a certified artefact that cannot be
retrained~\cite{guo2017calibration,kull2019beyond}. The literature that developed
these methods, however, is organised around a different target. Expected
calibration error (ECE) asks whether confidence matches accuracy \emph{on average
within confidence bands}; selective risk asks whether confidence \emph{orders} the
errors. The two coincide in the ideal case and come apart in practice, because a
monotone transform of the confidence score leaves every risk--coverage curve
unchanged while moving ECE (Proposition~\ref{prop:inv}).

This paper asks what is available to a post-hoc method if it optimises for
ordering rather than for average agreement, and the answer starts from what the
frozen model already computes. A linear head reduces a $d$-dimensional penultimate
representation to a handful of logits, discarding everything about the geometry of
that representation---which regions are dense, which sit near a decision boundary,
which are isolated. A layer that reads the latent vector directly can recover part
of that geometry and say something the softmax cannot. We build such a layer,
APEX, from three kernel experts of deliberately different inductive bias fitted on
a low-dimensional projection of the latent space, and mix their posteriors with
the original softmax. The construction is what is new: the components are
individually standard, but stacking them over a shared subspace and
gating the base model on its own confidence yields a score that ranks errors far
better than any of its parts. Our contributions are four. We give \textbf{a post-hoc fusion layer for frozen
backbones} that reduces AURC by 44.3\% and wins all 18 dataset $\times$ backbone
configurations against each of six baselines, at the smallest $p$-value the
Wilcoxon test can return at $n=18$. We establish \textbf{a separation between two
reliability criteria}, proving selective risk invariant under monotone
reparametrisation where ECE is not, and measuring the gap: in 15 of 18
configurations the lowest-ECE method is not the lowest-AURC one. We report our own
ECE plainly, including where a baseline beats it. We show \textbf{the objective governs calibration}---refitting the same
experts under a focal objective instead of the strictly proper one costs
$1.8\times10^{-2}$ ECE and buys nothing. And we work from \textbf{an audited
corpus}: we hash every image before use and report what we removed, 1{,}337 files
including 798 segmentation masks a file glob had been loading as ultrasound
scans.

\begin{figure*}[t]
\centering
\includegraphics[width=\textwidth]{fig1_architecture.pdf}
\caption{APEX in full. The backbone and its linear head are frozen; everything to
the right of the dashed boundary is fitted post-hoc on held-out splits and adds no
gradient step to the network. Latent vector and descriptor block are projected to a
shared 64-dimensional subspace by a truncated SVD, three kernel experts are
evaluated on it, and their posteriors are mixed with the base softmax under a gate
that discounts the base model exactly when the base model is unsure.}
\label{fig:arch}
\end{figure*}

\section{Background}

\subsection{Post-hoc calibration}
Temperature scaling divides the logits by one learned scalar and is the standard
reference point, being both effective and unable to change any
prediction~\cite{guo2017calibration}. Vector and matrix scaling generalise it to
per-class affine maps, isotonic regression fits a monotone transform per
class~\cite{zadrozny2002transforming}, and Dirichlet calibration places a
regularised log-linear model on the simplex~\cite{kull2019beyond}. Across the
family the map acts on the model's own output and is fitted to a criterion
evaluated on that output alone.

\subsection{Measuring calibration}
ECE bins predictions by confidence and averages the gap between confidence and
accuracy~\cite{naeini2015obtaining}. Binning makes the estimate optimistic---errors
of opposite sign cancel inside a bin---which motivates debiased
estimators~\cite{kumar2019verified} and class-wise
variants~\cite{nixon2019measuring}. Strictly proper scoring rules, Brier and
log-loss, avoid binning and are minimised only by the true conditional
distribution~\cite{brier1950verification,gneiting2007strictly}. We report all of
these and treat the proper scores as primary evidence on probability quality.

\subsection{Selective prediction and stacking}
Selective classification pairs a predictor with a rejection
rule~\cite{elyaniv2010foundations}. Sweeping the threshold traces risk against
coverage; the area under that curve (AURC) summarises how well the confidence
score ranks errors~\cite{geifman2017selective}; its excess form subtracts the
value an oracle ranking would attain. Learned reject options such as
SelectiveNet~\cite{geifman2019selectivenet} require training, which a frozen
backbone rules out; maximum softmax probability remains the standard training-free
baseline~\cite{hendrycks2017baseline}. APEX is a stacked
generalisation~\cite{wolpert1992stacked} whose base learners see a projection of
the frozen latent space rather than the raw image, and whose meta-model is a gated
convex combination fitted under a strictly proper rule.

\section{Method}

\subsection{Setting and notation}
Fig.~\ref{fig:arch} shows the construction end to end; this section defines it.
Let $f$ be a trained classifier over $K$ classes with fixed parameters. For an
image $X_i$ we read the penultimate representation
$\mathbf{z}_i \in \mathbb{R}^{d}$ and the softmax output
$\mathbf{p}^{\mathrm{b}}_i \in \Delta^{K-1}$, and compute a 62-dimensional
descriptor vector $\mathbf{m}_i$ (intensity statistics, grey-level co-occurrence
features, local binary patterns, shape moments, wavelet energies). Labelled data is
split three ways: $\mathcal{D}_{\mathrm{fit}}$ for the experts,
$\mathcal{D}_{\mathrm{cal}}$ for the mixture weights, and an evaluation fold no
component ever sees.

\subsection{A shared subspace for heterogeneous features}
Latent coordinates and descriptors live on incompatible scales, so we standardise
the concatenation $\mathbf{w}_i = [\mathbf{z}_i ; \mathbf{m}_i]$ and project it to
a common subspace,
\begin{equation}
\mathbf{u}_i \;=\; \Phi\!\left( \mathbf{w}_i \right),
\qquad \Phi : \mathbb{R}^{d+62} \rightarrow \mathbb{R}^{r},
\label{eq:proj}
\end{equation}
where $\Phi$ standardises each coordinate and then applies a truncated singular
value decomposition to $r = 64$ components, fitted on $\mathcal{D}_{\mathrm{fit}}$
alone. The projection is linear; every nonlinearity in APEX lives in the experts.
Compression is not the point: kernel machines degrade in high dimension, and
$r \ll d$ is what makes the RBF bandwidth meaningful.

\subsection{Three experts with different inductive biases}
On $\mathbf{u}$ we fit three probabilistic classifiers chosen so their errors are
unlikely to coincide. An RBF support vector machine responds to local density with
a smooth, globally regularised boundary; a cubic-polynomial SVM captures low-order
curvature the RBF kernel smooths away; a $k$-nearest-neighbour classifier is purely
local, and asks what a linear head cannot---are this sample's neighbours of one
class or of several? Both SVMs use $C=1$, with the scaled RBF bandwidth and degree
three respectively, and turn margins into posteriors by Platt
scaling~\cite{platt1999probabilistic}; the $k$-NN posterior is the
inverse-distance-weighted vote of the $k=5$ nearest neighbours. We write them $\mathbf{p}^{\mathrm{rbf}}_i$,
$\mathbf{p}^{\mathrm{poly}}_i$, $\mathbf{p}^{\mathrm{knn}}_i$.

\subsection{Fusion and the confidence gate}
The fused posterior is a convex combination of the base softmax and the three
experts,
\begin{equation}
\hat{\mathbf{p}}^{\,\circ}_i \;=\; w_1\,\mathbf{p}^{\mathrm{b}}_i \;+\; w_2\,\mathbf{p}^{\mathrm{rbf}}_i
\;+\; w_3\,\mathbf{p}^{\mathrm{poly}}_i \;+\; w_4\,\mathbf{p}^{\mathrm{knn}}_i,
\label{eq:fuse}
\end{equation}
with $w_j \ge 0$, $\sum_j w_j = 1$. Fixed weights trust the base model equally on
the cases it understands and the cases it does not, so we make the first weight a
function of its own confidence. Writing $c_i = \max_k p^{\mathrm{b}}_{ik}$ and
fixing a threshold $\tau$,
\begin{equation}
\tilde{w}_1(c_i) \;=\;
\begin{cases}
w_1, & c_i \ge \tau,\\[2pt]
\lambda\,w_1, & c_i < \tau,
\end{cases}
\qquad 0 < \lambda < 1,
\label{eq:gate}
\end{equation}
and the mass $\left(1-\lambda\right) w_1$ freed below the threshold is
redistributed uniformly across the experts. The gated predictor
$\hat{\mathbf{p}}_i$ is Eq.~\ref{eq:fuse} with $w_1$ replaced by
$\tilde{w}_1(c_i)$, and stays on the simplex. The gate is crude on purpose: it does
not work out which expert is right, it withdraws trust from the base model where
the base model has already said it has none. We use $\tau = 0.85$,
$\lambda = 0.20$ throughout; Section~\ref{sec:ablation} shows the result is
insensitive to $\tau$.

\subsection{Fitting the weights}
The weights are obtained on $\mathcal{D}_{\mathrm{cal}}$ by minimising negative
log-likelihood,
\begin{equation}
\mathbf{w}^{\star} \;=\; \arg\min_{\mathbf{w} \in \Delta^{3}}
\; -\frac{1}{|\mathcal{D}_{\mathrm{cal}}|}
\sum_{i \in \mathcal{D}_{\mathrm{cal}}} \log \hat{p}^{\,\circ}_{i, y_i},
\label{eq:nll}
\end{equation}
a convex problem in $\mathbf{w}$. The objective is the \emph{ungated} mixture:
keeping the threshold out of it preserves convexity, and the gate, which only
redistributes a fixed fraction of $w_1$, is applied at inference with the weights
so obtained. Negative log-likelihood is a strictly proper scoring rule, so its
minimiser over a well-specified family is the true conditional
distribution~\cite{gneiting2007strictly}. Objectives built to sharpen decisions
rather than to score probabilities---focal loss being the common
example~\cite{mukhoti2020calibrating}---carry no such guarantee, and
Section~\ref{sec:ablation} measures what substituting one costs. APEX adds no
gradient step to the backbone: fitting takes under a minute per configuration on
CPU for our largest dataset, and inference adds three kernel evaluations and a
weighted sum.

\section{Experimental protocol}

\subsection{Data}
Six public corpora, one per modality, 32{,}251 images in total after the audit
below: blood smear microscopy (12{,}513 images, 4
classes)~\cite{mooney2018blood}, brain MRI (7{,}013, 4)~\cite{nickparvar2021brain},
chest radiography (3{,}994, 2)~\cite{kermany2018identifying}, endoscopy (3{,}990,
4)~\cite{heartzhacker2023medical}, retinal OCT (3{,}963,
4)~\cite{kermany2018identifying} and breast ultrasound
(778, 3)~\cite{aldhabyani2020dataset}. Every image carries a matched descriptor
row on all six.

\subsection{Data hygiene}
\label{sec:hygiene}
Duplicate and mislabelled files inflate cross-validated results silently: copies
of one image land in different folds and the model is tested on material it
memorised~\cite{barz2020purging}. We hashed all 33{,}588 candidate files and
quarantined 1{,}337 (4.0\%): 798 binary segmentation masks a file-extension glob
had been loading as ultrasound scans, 531 byte-identical copies within a class,
and both members of the four groups in which one image appeared under two labels. The masks were the most consequential: they halved the
nominal ultrasound corpus and made its classes separable by silhouette alone. Every
number here is computed after that removal.

\subsection{Backbones and heads}
We use three architecture families of different design: ViT-S/16 ($d=384$)~\cite{dosovitskiy2020image},
ConvNeXt-T ($d=768$)~\cite{liu2022convnet} and MobileNetV3-L ($d=1280$)~\cite{howard2019searching}.
Every backbone parameter is held fixed at its ImageNet value and only a linear
head with dropout is trained, ten epochs per fold, the checkpoint selected on an
inner split so the evaluation fold is never used for model selection. One caveat,
stated precisely: MobileNetV3 uses batch normalisation, so its running statistics
adapt to the training split---never to an evaluation fold---while its weights stay
frozen. ViT and ConvNeXt normalise per layer and are unaffected.

\subsection{Protocol and metrics}
Every dataset $\times$ backbone pair is evaluated under stratified five-fold
cross-validation, giving 90 trained heads and 18 paired configurations; within each
fold the training portion splits 80/20 into $\mathcal{D}_{\mathrm{fit}}$ and
$\mathcal{D}_{\mathrm{cal}}$; all fifteen metrics appear in
Table~\ref{tab:allmetrics}. Isotonic regression is a step function and ties 71\%
of its confidences, so we report AURC as the expectation over tie-breaking, which
is independent of input order and coincides with the plain definition when
confidences are distinct. Baselines are the uncalibrated softmax, MC dropout
with 30 stochastic passes~\cite{gal2016dropout}, temperature and vector
scaling~\cite{guo2017calibration}, isotonic
regression~\cite{zadrozny2002transforming} and Dirichlet
calibration~\cite{kull2019beyond}. Following standard practice for multi-dataset
comparison~\cite{demsar2006statistical}, configurations are the unit of analysis
and methods are compared with the two-sided Wilcoxon signed-rank
test~\cite{wilcoxon1945individual}. MC dropout's stochasticity was verified rather
than assumed: the passes differed in all 90 folds.

\begin{table*}[t]
\centering
\caption{Mean $\pm$ standard deviation over the 18 dataset$\,\times\,$backbone configurations. The last two rows count the configurations APEX wins and give the worst two-sided Wilcoxon $p$ against any single baseline; $7.6\times10^{-6}$ is the smallest value the test can return at $n{=}18$.}
\label{tab:headline}
\setlength{\tabcolsep}{4.0pt}\renewcommand{\arraystretch}{0.86}
\footnotesize
\begin{tabular}{l ccccccc}
\toprule
Method & BAcc $\uparrow$ & ECE $\downarrow$ & SCE $\downarrow$ & Brier $\downarrow$ & NLL $\downarrow$ & AURC $\downarrow$ & Risk@90 $\downarrow$ \\
\midrule
Uncalibrated softmax & 0.8523\,\tiny{$\pm$0.0963} & 0.0495\,\tiny{$\pm$0.0338} & 0.0468\,\tiny{$\pm$0.0311} & 0.2062\,\tiny{$\pm$0.1169} & 0.3710\,\tiny{$\pm$0.1994} & 0.0486\,\tiny{$\pm$0.0510} & 0.1071\,\tiny{$\pm$0.0860} \\
MC dropout \cite{gal2016dropout} & 0.8531\,\tiny{$\pm$0.0942} & 0.0571\,\tiny{$\pm$0.0337} & 0.0497\,\tiny{$\pm$0.0297} & 0.2088\,\tiny{$\pm$0.1166} & 0.3773\,\tiny{$\pm$0.1987} & 0.0488\,\tiny{$\pm$0.0509} & 0.1078\,\tiny{$\pm$0.0859} \\
Temperature scaling \cite{guo2017calibration} & 0.8523\,\tiny{$\pm$0.0963} & 0.0328\,\tiny{$\pm$0.0278} & 0.0393\,\tiny{$\pm$0.0289} & 0.2015\,\tiny{$\pm$0.1135} & 0.3597\,\tiny{$\pm$0.1931} & 0.0484\,\tiny{$\pm$0.0510} & 0.1066\,\tiny{$\pm$0.0859} \\
Vector scaling \cite{guo2017calibration} & 0.8583\,\tiny{$\pm$0.0934} & 0.0311\,\tiny{$\pm$0.0244} & 0.0313\,\tiny{$\pm$0.0252} & 0.1931\,\tiny{$\pm$0.1083} & 0.3488\,\tiny{$\pm$0.1821} & 0.0444\,\tiny{$\pm$0.0472} & 0.1000\,\tiny{$\pm$0.0808} \\
Isotonic regression \cite{zadrozny2002transforming} & 0.8569\,\tiny{$\pm$0.0941} & 0.0347\,\tiny{$\pm$0.0248} & 0.0328\,\tiny{$\pm$0.0265} & 0.1968\,\tiny{$\pm$0.1113} & 0.4466\,\tiny{$\pm$0.2864} & 0.0484\,\tiny{$\pm$0.0500} & 0.1018\,\tiny{$\pm$0.0829} \\
Dirichlet calibration \cite{kull2019beyond} & 0.8580\,\tiny{$\pm$0.0960} & \textbf{0.0299\,\tiny{$\pm$0.0228}} & 0.0304\,\tiny{$\pm$0.0249} & 0.1930\,\tiny{$\pm$0.1089} & 0.3443\,\tiny{$\pm$0.1853} & 0.0447\,\tiny{$\pm$0.0477} & 0.1014\,\tiny{$\pm$0.0829} \\
\midrule
\textbf{APEX (ours)} & \textbf{0.8941\,\tiny{$\pm$0.0854}} & 0.0339\,\tiny{$\pm$0.0206} & \textbf{0.0297\,\tiny{$\pm$0.0225}} & \textbf{0.1432\,\tiny{$\pm$0.0956}} & \textbf{0.2589\,\tiny{$\pm$0.1635}} & \textbf{0.0271\,\tiny{$\pm$0.0329}} & \textbf{0.0652\,\tiny{$\pm$0.0678}} \\
\midrule
APEX better in & 17--18/18 & 6--16/18 & 9--18/18 & 18/18 & 18/18 & 18/18 & 17--18/18 \\
worst $p$ vs.\ any row & $<10^{-4}$ & $0.87$ & $0.80$ & $<10^{-4}$ & $<10^{-4}$ & $<10^{-4}$ & $<10^{-3}$ \\
\bottomrule
\end{tabular}
\end{table*}


\section{Results}

\subsection{Selective risk}
Averaged over the 18 configurations (Table~\ref{tab:headline}), APEX reduces AURC
from $4.86\times10^{-2}$ to $2.71\times10^{-2}$---44.3\% against the uncalibrated
softmax, 39.0\% against the strongest baseline---and risk at 90\% coverage from
10.71\% to 6.52\%. The margin is not carried by a subset of easy cases.
Fig.~\ref{fig:rc}(a) shows the mean risk--coverage curve below every baseline at
every coverage, the gap growing from 0.7 points of risk at 30\% coverage to 3.5
points at 90\%---the range a triage system would use.
Fig.~\ref{fig:rc}(b) makes the comparison as strict as it can be made---APEX
against whichever baseline was best \emph{for that configuration}---and all
eighteen points fall below the diagonal. Fig.~\ref{fig:rc}(c) shows the effect in
every modality, from 18\% on endoscopy to 82\% on blood smears.

The proper scores move with it: Brier falls 30.6\% and NLL 30.2\%, each 18/18
against all six baselines at $p = 7.6\times10^{-6}$. Balanced accuracy rises 4.18
points, 18/18 against five of the six and 17/18 against Dirichlet
($p = 2.3\times10^{-5}$). Because APEX may change the arg-max this is a genuine
accuracy gain rather than a rescaling---obtained with the backbone frozen.

\begin{figure*}[t]
\centering
\includegraphics[width=\textwidth]{fig2_selective.pdf}
\caption{Selective performance. (a) Mean risk--coverage curve over the 18
configurations, traced from the per-sample predictions of all 90 folds; the band
is one standard error. (b) Per-configuration AURC against the best baseline
\emph{for that configuration}. (c) AURC per modality, averaged over the three
backbones. Panels (b) and (c) use the same per-fold values as
Table~\ref{tab:headline}.}
\label{fig:rc}
\end{figure*}

\subsection{Calibration is a different axis}
ECE does not follow. APEX cuts it from $4.95\times10^{-2}$ to
$3.39\times10^{-2}$, but the calibrators land in the same place---$3.28$, $3.11$
and $3.47\times10^{-2}$ for temperature, vector and isotonic scaling---and the
paired tests separate none of them from APEX ($p = 0.37$, $0.13$, $0.87$).
Dirichlet calibration reaches $2.99\times10^{-2}$ and beats APEX in 12 of the 18
configurations ($p = 0.030$). APEX also has the highest maximum calibration error
of anything we ran. Fig.~\ref{fig:cal}(a) shows why: APEX sits above the diagonal
across the whole range, predicting less confidence than it earns, while Dirichlet
tracks it closely. Under-confidence is the direction a triage system can live with, but ECE charges
for it either way.
What APEX gets right is the global level: its average confidence is within $0.9$
points of its accuracy, closer than any other method. The ECE it pays comes from
spread inside the bins, not from being systematically too sure of itself.

The divergence has a structural reason that does not depend on our data.

\begin{proposition}
\label{prop:inv}
Let $g$ be strictly increasing on $[0,1]$. Replacing every confidence $c_i$ by
$g(c_i)$ leaves the risk--coverage curve unchanged, and with it AURC, E-AURC and
selective risk at every coverage. Expected calibration error has no such
invariance.
\end{proposition}

\begin{IEEEproof}
At coverage $\kappa$ the curve reports the empirical error over the
$\lceil \kappa n\rceil$ samples of largest confidence, so it depends on the
$c_i$ only through the ordering they induce, which a strictly increasing $g$
preserves. ECE compares confidence \emph{values} with accuracies,
$\sum_b (n_b/n)|\bar c_b - \bar a_b|$, so it moves when those values move.
\end{IEEEproof}

The prediction is exact and it holds numerically: applying $c^2$, $\sqrt{c}$,
$\sigma(6(c-\tfrac12))$ and $\tfrac12+\tfrac12c$ to the uncalibrated confidences
of all 18 configurations leaves AURC identical in every digit, largest deviation
$0$, while ECE ranges from $4.95\times10^{-2}$ to $1.44\times10^{-1}$. A
rescaling that cannot change which cases a system defers moves the calibration
error it reports by a factor of $2.9$.

The consequence shows up in the rankings, and Fig.~\ref{fig:cal}(b)
quantifies it. Across the 18 configurations APEX holds mean rank 1.00
on AURC, on Brier and on NLL---it is first in every single configuration on all
three---and mean rank 3.78 on ECE, behind Dirichlet, vector scaling and temperature
scaling. Choosing
a method by calibration error would have selected the wrong method for deferral in
15 of the 18 configurations. A post-hoc layer should be evaluated against the
criterion it will be deployed under, and for abstention that criterion is the
ordering, not the average gap.

\begin{figure*}[t]
\centering
\includegraphics[width=\textwidth]{fig3_calibration.pdf}
\caption{Calibration and ordering are not the same property. (a) Reliability with
equal-mass bins formed inside each configuration and averaged over the 18, the
same way the ECE column of Table~\ref{tab:headline} is computed; band is one
standard error. (b) Mean rank over the 18 configurations on ECE and on AURC; APEX
is first on AURC in every configuration and mid-field on ECE, and the other
methods reorder between the two axes as well.}
\label{fig:cal}
\end{figure*}

\subsection{What each component contributes}
\label{sec:ablation}
Fig.~\ref{fig:abl}(a) removes one part at a time and re-fits. The RBF expert is
load-bearing: without it AURC rises by $2.9\times10^{-3}$ and ECE by
$9.0\times10^{-3}$ ($p<10^{-3}$ for both). The $k$-NN expert is second for ordering
($+1.5\times10^{-3}$ AURC, $p=2.5\times10^{-4}$) and costs a little calibration
error to include---a component that helps the ranking while mildly hurting the
average gap, the same asymmetry seen at the method level. The cubic SVM and the
descriptor block each give a small but significant AURC improvement
($p = 1.6\times10^{-2}$ and $1.6\times10^{-3}$). The gate's contribution is
specific: it improves ECE reliably ($p=8.4\times10^{-4}$) but leaves AURC
unchanged ($p = 0.20$).

No expert reproduces the fused result alone: the RBF SVM, the strongest, reaches
AURC $3.12\times10^{-2}$ against $2.71\times10^{-2}$ for the fusion
($p = 1.5\times10^{-5}$), and an equal-weight blend without gating or fitted
weights reaches $2.85\times10^{-2}$.

The largest single effect is the objective. Refitting the same experts and gate
under a focal loss instead of Eq.~\ref{eq:nll} raises ECE by $1.8\times10^{-2}$
and NLL by $1.3\times10^{-2}$ ($p = 7.6\times10^{-6}$ for both) while leaving
AURC and accuracy statistically unchanged: the fused probabilities are only as
honest as the rule used to fit them. The gate threshold is not tuned---sweeping
$\tau$ over $[0.60,0.95]$ moves AURC by 0.5\% of its own value
(Fig.~\ref{fig:abl}(b)).

\begin{figure}[t]
\centering
\includegraphics[width=\columnwidth]{fig4_ablation.pdf}
\caption{(a) Change in AURC and ECE when one component is removed; bars to the
right are degradations. Wilcoxon: $^{**}p<0.01$, $^{*}p<0.05$. (b) $\tau$ is not a
tuned parameter---AURC varies by 0.5\% of its own value across the sweep.}
\label{fig:abl}
\end{figure}

\subsection{Full metric detail}
Table~\ref{tab:allmetrics} reports all fifteen metrics per backbone. Two patterns
hold across architectures: the selective and proper scores improve for ViT,
ConvNeXt and MobileNetV3 alike, with the largest absolute gains on the weakest
backbone, while maximum calibration error and over-confidence error are the two
metrics on which APEX trails even the uncalibrated model on every backbone.

\begin{table*}[t]
\centering
\caption{All fifteen metrics per backbone, averaged over the six datasets (each entry is itself a mean over five folds). \emph{None} is the uncalibrated softmax; Dirichlet is the strongest baseline by mean rank on ECE, Brier and NLL. Arrows give the preferred direction; best of each block in bold.}
\label{tab:allmetrics}
\setlength{\tabcolsep}{3.0pt}\renewcommand{\arraystretch}{0.74}
\scriptsize
\begin{tabular}{l c *{3}{ccc}}
\toprule
 & & \multicolumn{3}{c}{ViT-S/16 ($d{=}384$)} & \multicolumn{3}{c}{ConvNeXt-T ($d{=}768$)} & \multicolumn{3}{c}{MobileNetV3-L ($d{=}1280$)} \\
\cmidrule(lr){3-5}\cmidrule(lr){6-8}\cmidrule(lr){9-11}
Metric & Dir. & None & Dirich. & \textbf{APEX} & None & Dirich. & \textbf{APEX} & None & Dirich. & \textbf{APEX} \\
\midrule
Balanced accuracy & $\uparrow$ & 0.8620 & 0.8653 & \textbf{0.9040} & 0.8621 & 0.8777 & \textbf{0.8961} & 0.8329 & 0.8310 & \textbf{0.8822} \\
Macro F1 & $\uparrow$ & 0.8630 & 0.8679 & \textbf{0.9073} & 0.8632 & 0.8751 & \textbf{0.8981} & 0.8280 & 0.8311 & \textbf{0.8839} \\
Macro AUROC & $\uparrow$ & 0.9606 & 0.9617 & \textbf{0.9769} & 0.9656 & 0.9659 & \textbf{0.9769} & 0.9496 & 0.9489 & \textbf{0.9692} \\
ECE & $\downarrow$ & 0.0394 & \textbf{0.0297} & 0.0333 & 0.0455 & \textbf{0.0299} & 0.0329 & 0.0635 & \textbf{0.0301} & 0.0355 \\
Debiased ECE & $\downarrow$ & 0.0150 & \textbf{0.0072} & 0.0127 & 0.0212 & \textbf{0.0070} & 0.0126 & 0.0362 & \textbf{0.0078} & 0.0143 \\
Class-wise SCE & $\downarrow$ & 0.0409 & 0.0290 & \textbf{0.0279} & 0.0440 & \textbf{0.0293} & 0.0301 & 0.0555 & 0.0329 & \textbf{0.0312} \\
Maximum CE & $\downarrow$ & \textbf{0.3329} & 0.3588 & 0.3769 & 0.3738 & \textbf{0.2894} & 0.3808 & \textbf{0.3030} & 0.3244 & 0.3464 \\
Over-confidence error & $\downarrow$ & \textbf{0.6646} & 0.7071 & 0.6842 & \textbf{0.6561} & 0.7083 & 0.6874 & \textbf{0.6056} & 0.6749 & 0.6783 \\
Brier score & $\downarrow$ & 0.1938 & 0.1833 & \textbf{0.1269} & 0.1858 & 0.1734 & \textbf{0.1385} & 0.2389 & 0.2224 & \textbf{0.1641} \\
Brier skill score & $\uparrow$ & 0.7208 & 0.7356 & \textbf{0.8142} & 0.7323 & 0.7497 & \textbf{0.7987} & 0.6561 & 0.6799 & \textbf{0.7619} \\
NLL & $\downarrow$ & 0.3473 & 0.3290 & \textbf{0.2311} & 0.3327 & 0.3081 & \textbf{0.2495} & 0.4330 & 0.3957 & \textbf{0.2962} \\
AURC & $\downarrow$ & 0.0463 & 0.0414 & \textbf{0.0244} & 0.0390 & 0.0360 & \textbf{0.0243} & 0.0606 & 0.0566 & \textbf{0.0326} \\
Excess AURC & $\downarrow$ & 0.0333 & 0.0295 & \textbf{0.0178} & 0.0265 & 0.0252 & \textbf{0.0168} & 0.0404 & 0.0376 & \textbf{0.0233} \\
Risk at 90\% coverage & $\downarrow$ & 0.0993 & 0.0923 & \textbf{0.0558} & 0.0938 & 0.0872 & \textbf{0.0614} & 0.1284 & 0.1248 & \textbf{0.0784} \\
Risk at 80\% coverage & $\downarrow$ & 0.0748 & 0.0688 & \textbf{0.0396} & 0.0699 & 0.0653 & \textbf{0.0433} & 0.1025 & 0.0971 & \textbf{0.0586} \\
\bottomrule
\end{tabular}
\end{table*}


\section{Discussion}

The reason is simple. A frozen linear head is a projection, and projections lose
information: two samples with nearly the same logits can sit
in very different neighbourhoods of the latent space---one deep inside a class
cluster, the other on a boundary---and the softmax cannot express the difference.
The kernel experts can, an RBF machine reading local density and a $k$-NN
classifier neighbourhood purity, so fusing their posteriors with the softmax
reinstates a signal about \emph{which} predictions are unsafe.

It also explains the asymmetry. A parametric calibrator does the opposite of what
APEX does---temperature scaling provably cannot change any prediction, so it moves
ECE while leaving the ordering untouched. Proposition~\ref{prop:inv} is the formal
statement of that separation, and it is why the two families are complementary:
composing a scaling layer on a fusion layer is a natural next step we have not
tested.

The numbers translate directly: at 90\% coverage selective risk falls from 10.71\%
to 6.52\%, so roughly two in five of the errors that would have reached a clinician
are routed for review instead, at no change to the network.

\section{Limitations}
\label{sec:limits}
Four limitations bound these claims. The backbones are frozen and only linear heads
trained; we have not shown the gains survive full fine-tuning, where a stronger
head may already encode some of the geometry the experts recover. That is the most
important experiment we did not run.

Every result is from a single seed---fold assignment, the fit/calibration split
and head initialisation are fixed---so the 18-of-18 margins carry unmeasured seed
variance. The breast ultrasound corpus is small (778 images after the audit) and
drives the widest fold-to-fold spread; all six corpora are public, with no
external or prospective validation.

Finally, APEX does not improve expected calibration error over the best parametric
baselines and is worse than one of them. A deployment needing a calibrated
probability should use a calibrator, or compose one with APEX; our claim is
confined to the ordering of confidence and the abstention decisions that depend on
it.

\section{Conclusion}

We asked what a post-hoc layer can do for a frozen medical image classifier built
for abstention rather than average calibration. Reading the latent representation
the classifier discards, fitting three kernel experts over a shared projection of
it and mixing their posteriors with the softmax under a confidence-conditional
gate reduces AURC by 44.3\% and wins all 18 configurations against each of six
baselines, with Brier and log-likelihood doing the same and balanced accuracy
winning at least 17 of 18. Expected calibration error does not improve beyond what
parametric calibrators achieve---the consequence of optimising a different
property, not a shortfall to explain away. Selective risk is invariant to monotone
rescaling where ECE is not, and the two disagree about which method to deploy in
15 of 18 cases; the criterion should be chosen before the method is.

\bibliographystyle{IEEEtran}
\bibliography{references}

\end{document}
